\documentclass[11pt,a4paper]{article}

\usepackage[utf8]{inputenc}
\usepackage[T1]{fontenc}
\usepackage[italian]{babel}
\usepackage{amsmath,amssymb,amsfonts}
\usepackage{geometry}
\usepackage{booktabs}
\usepackage{array}
\usepackage{hyperref}
\usepackage{listings}
\usepackage{xcolor}
\usepackage{longtable}

\geometry{margin=2.5cm}
\hypersetup{colorlinks=true, linkcolor=black, urlcolor=blue, citecolor=black}

\lstdefinestyle{code}{
  basicstyle=\ttfamily\small,
  breaklines=true,
  frame=single,
  columns=fullflexible,
  backgroundcolor=\color{gray!5}
}
\lstset{style=code}

\title{SemaGraph v0.6\\
\large Kernel deterministico per ancoraggio di stati osservati e compressione di catene di transizione in $Q_6$}
\author{Bozza tecnica per sviluppo e validazione}
\date{Giugno 2026}

\begin{document}
\maketitle

\begin{abstract}
SemaGraph v0.6 propone un kernel deterministico per rappresentare cambiamenti osservati come catene finite di transizioni in uno spazio booleano a sei dimensioni, indicato come $Q_6$. Ogni stato è codificato da una parola binaria di sei bit; ogni transizione diretta tra due stati è rappresentata da una maschera XOR; ogni catena di trasformazione è comprimibile in una firma operativa costituita da stato sorgente, stato target, maschere ordinate, mutazione netta e distanza cumulativa. Il sistema non usa modelli linguistici per inferire stati o parametri: l'ancoraggio iniziale deriva da osservazioni fisiche, sensori, soglie o parametrizzazioni dichiarate. I modelli linguistici possono intervenire soltanto a valle, per sintetizzare policy o spiegazioni su firme già calcolate. L'obiettivo è ridurre inferenza ricorrente, drift interpretativo e costo computazionale nella lettura di traiettorie dinamiche, mantenendo tracciabilità verso i dati osservati.
\end{abstract}

\section{Premessa}
Molti sistemi complessi producono sequenze di stati più che singoli stati isolati. Processi organizzativi, modelli finanziari, fenomeni biologici, eventi fisici discretizzati e simulazioni stocastiche possono essere osservati come traiettorie: una configurazione iniziale cambia, attraversa regimi intermedi e raggiunge una configurazione finale. Il problema non è soltanto classificare lo stato corrente, ma rappresentare come il sistema vi è arrivato.

SemaGraph v0.6 nasce per formalizzare questo passaggio. Il kernel non pretende di sostituire modelli fisici, chimici, biologici, statistici o linguistici. Esso opera come layer di compressione transizionale: prende stati già osservati o parametrizzati, calcola transizioni deterministiche, comprime catene e genera firme confrontabili.

La tesi principale è la seguente: se gli stati sono ancorati deterministicamente a osservazioni o parametri noti, allora la catena di trasformazione può essere rappresentata come sequenza di transizioni finite, replayable e comprimibili. L'inferenza viene confinata fuori dal nucleo transizionale.

\section{Modello formale}
\subsection{Spazio degli stati}
Sia $B = \{0,1\}$. Lo spazio State64 è definito come:
\begin{equation}
Q_6 = B^6.
\end{equation}
Ogni stato $s \in Q_6$ è una tupla ordinata di sei bit:
\begin{equation}
s = (b_1,b_2,b_3,b_4,b_5,b_6), \quad b_i \in \{0,1\}.
\end{equation}
Lo spazio contiene esattamente:
\begin{equation}
|Q_6| = 2^6 = 64
\end{equation}
stati possibili. In implementazione, ogni stato può essere rappresentato come intero unsigned a 8 bit usando soltanto i sei bit meno significativi o come stringa canonica di sei caratteri binari.

\subsection{Transizione elementare}
Date due configurazioni $s,t \in Q_6$, la transizione diretta da $s$ a $t$ è definita dalla maschera:
\begin{equation}
m = s \oplus t,
\end{equation}
dove $\oplus$ indica XOR bit-a-bit. La maschera $m$ specifica quali posizioni cambiano. Se $m_i = 1$, il bit $i$ è mutato; se $m_i = 0$, il bit $i$ è conservato.

La distanza di trasformazione è la distanza di Hamming:
\begin{equation}
d(s,t) = \sum_{i=1}^{6} |s_i - t_i| = \operatorname{popcount}(s \oplus t).
\end{equation}
Questa misura non interpreta il cambiamento: misura quante coordinate dello stato sono mutate.

\subsection{Matrice completa delle transizioni}
Poiché $Q_6$ contiene 64 stati, il numero di coppie ordinate stato-stato è:
\begin{equation}
64 \times 64 = 4096.
\end{equation}
Queste 4096 coppie costituiscono la base completa delle transizioni dirette, inclusa l'identità $s \rightarrow s$. Le transizioni trasformative non identitarie sono:
\begin{equation}
64 \times 63 = 4032.
\end{equation}
L'indice computazionale di una transizione può essere definito come:
\begin{equation}
I(s,t) = 64 \cdot \operatorname{idx}(s) + \operatorname{idx}(t),
\end{equation}
dove $\operatorname{idx}(s) \in [0,63]$ è la rappresentazione intera dello stato sorgente.

\subsection{Modularità 3+3}
Ogni stato può essere scomposto in due moduli da tre bit:
\begin{equation}
s = (l,u), \quad l \in B^3,\ u \in B^3.
\end{equation}
La scomposizione 3+3 non introduce significato ereditato; è una proprietà combinatoria utile per classificare se una transizione rimane interna a un modulo o attraversa entrambi i moduli. Date le distanze:
\begin{equation}
d_l = d(l_s,l_t), \quad d_u = d(u_s,u_t),
\end{equation}
possiamo classificare il campo di mutazione come: nessun cambiamento, solo modulo inferiore, solo modulo superiore o entrambi i moduli.

\section{Ancoraggio deterministico degli stati}
\subsection{Osservazione e parametrizzazione}
SemaGraph v0.6 assume che lo stato non venga inferito linguisticamente. Lo stato è prodotto da regole di ancoraggio esplicite applicate a parametri osservati:
\begin{equation}
A: X \rightarrow Q_6,
\end{equation}
dove $X$ è lo spazio dei vettori osservativi o parametrici. Ogni bit $b_i$ è calcolato da una regola dichiarata:
\begin{equation}
b_i = r_i(x), \quad r_i: X \rightarrow \{0,1\}.
\end{equation}
Una regola può essere una soglia, una condizione booleana, una categoria nota o un predicato su un parametro.

\subsection{Replayability}
Dato lo stesso vettore osservato $x$ e la stessa versione delle regole $R = (r_1, \ldots, r_6)$, il sistema produce sempre lo stesso stato:
\begin{equation}
A_R(x) = s.
\end{equation}
La determinazione dello stato è quindi replayable. L'incertezza di misura non viene negata: deve essere registrata come tolleranza, unità, fonte, timestamp e versione della regola. Tuttavia, questa incertezza appartiene al livello di misura, non al kernel di transizione.

\subsection{Boundary del modello linguistico}
Il modello linguistico non assegna stati, non modifica parametri e non inferisce misure mancanti. Può intervenire a valle:
\begin{itemize}
  \item per spiegare una firma transizionale;
  \item per proporre policy operative;
  \item per sintetizzare implicazioni leggibili;
  \item per confrontare policy già approvate.
\end{itemize}
Questa separazione impedisce che l'inferenza linguistica si propaghi dentro la catena deterministica.

\section{Compressione delle catene di transizione}
\subsection{Catena ordinata}
Una traiettoria osservata è una sequenza ordinata:
\begin{equation}
C = (s_0,s_1,\ldots,s_n), \quad s_i \in Q_6.
\end{equation}
Ogni passaggio produce una maschera:
\begin{equation}
m_i = s_{i-1} \oplus s_i.
\end{equation}
La catena può quindi essere rappresentata come:
\begin{equation}
C_M = (m_1,m_2,\ldots,m_n).
\end{equation}

\subsection{Mutazione netta e distanza cumulativa}
La mutazione netta è:
\begin{equation}
M_{net} = m_1 \oplus m_2 \oplus \cdots \oplus m_n = s_0 \oplus s_n.
\end{equation}
La distanza cumulativa è:
\begin{equation}
D_C = \sum_{i=1}^{n} \operatorname{popcount}(m_i).
\end{equation}
La mutazione netta esprime l'effetto complessivo; la sequenza ordinata conserva il percorso. Questi due livelli devono rimanere separati, perché catene diverse possono avere la stessa mutazione netta.

\subsection{Firma operativa}
Una firma transizionale minima può essere definita come:
\begin{equation}
\Sigma(C) = (s_0, s_n, (m_1,\ldots,m_n), M_{net}, D_C).
\end{equation}
Questa firma non sostituisce il dato grezzo, ma lo indicizza. Il dato originale resta necessario per audit, verifica e revisione.

\section{Architettura computazionale}
\subsection{Livelli}
L'architettura v0.6 distingue quattro livelli:
\begin{enumerate}
  \item osservazione e misura;
  \item ancoraggio deterministico a State64;
  \item processore di transizioni e compressione;
  \item sintesi di policy a valle.
\end{enumerate}
Il terzo livello è il nucleo computazionale. Esso può essere implementato con operazioni CPU elementari: XOR, AND, SHIFT, popcount e lookup table.

\subsection{Core Rust}
La versione v0.6 introduce un core Rust di fattibilità. La rappresentazione minima è:
\begin{lstlisting}[language=C]
state: u8          // valori ammessi 0..63
mutation_mask: u8  // valori ammessi 0..63
index: u16         // valori ammessi 0..4095
\end{lstlisting}
Le operazioni base sono:
\begin{lstlisting}[language=C]
next_state = state ^ mutation_mask
mutation   = source ^ target
distance   = popcount(source ^ target)
index      = source * 64 + target
\end{lstlisting}
Rust è adatto perché permette controllo numerico, portabilità, test di parità, compilazione verso WebAssembly e interfacce C ABI senza introdurre gestione manuale fragile della memoria.

\subsection{TypeScript come reference layer}
TypeScript rimane il livello di riferimento per specifica, API, esempi, JSON schema, integrazione web e strumenti AI. Rust non sostituisce immediatamente TypeScript: viene introdotto come candidato a diventare autorità computazionale dopo test di parità.

\section{Tool AI e sintesi di policy}
La versione v0.6 include un tool adapter per GPT/Claude. Il suo principio operativo è semplice:
\begin{quote}
Il modello può usare il kernel per processare stati e catene, ma non può alterare stati osservati, maschere o firme.
\end{quote}
Le funzioni principali sono:
\begin{itemize}
  \item \texttt{semagraph\_anchor\_state64}: ancoraggio di sei osservazioni note;
  \item \texttt{semagraph\_lookup\_transition64}: lookup di una transizione;
  \item \texttt{semagraph\_compress\_chain64}: compressione di una catena ordinata;
  \item \texttt{semagraph\_policy\_context64}: costruzione del contesto per policy synthesis.
\end{itemize}
Il modello linguistico opera quindi come redattore di policy o spiegatore semantico, non come motore di verità osservativa.

\section{Confronto con famiglie note}
SemaGraph non è un motore fisico, non è un planner general-purpose, non è un motore BPMN e non è un rule engine generico. Ha parentele con automi finiti, statecharts, modelli di transizione, rule engines e policy-as-code, ma il suo scopo è più ristretto: ancorare stati osservati e comprimere catene di trasformazione in firme operative.

\begin{longtable}{p{4cm}p{5cm}p{5cm}}
\toprule
Famiglia & Somiglianza & Differenza \\
\midrule
Automi finiti & Stati e transizioni finite & SemaGraph è orientato alla compressione di traiettorie osservate \\
Statecharts & Transizioni controllate e guardie & Non orchestra comportamento applicativo o UI \\
Rule engines & Predicati e regole & Non produce inferenza generale; valida transizioni e policy context \\
Policy-as-code & Regole operative formalizzabili & La policy è a valle della firma transizionale \\
Modelli stocastici & Traiettorie e percorsi & Non sceglie probabilisticamente il percorso; rappresenta deterministicamente i passaggi osservati \\
LLM toolchains & Output strutturato e validazione & LLM non è autorità su stati o transizioni \\
\bottomrule
\end{longtable}

\section{Ipotesi di validazione}
La validazione non deve dimostrare che SemaGraph ``descrive il mondo'' in senso generale. Deve testare ipotesi operative:
\begin{enumerate}
  \item riduce il tempo di classificazione di traiettorie osservate;
  \item riduce chiamate LLM e token necessari per produrre policy coerenti;
  \item aumenta replayability e auditabilità;
  \item conserva informazione decisionale sufficiente rispetto alla traiettoria grezza;
  \item evita falsi equivalenti prodotti da compressione eccessiva.
\end{enumerate}
Una prima validazione può usare traiettorie sintetiche e confrontare tre modalità: analisi libera, analisi LLM-only e analisi SemaGraph con policy synthesis a valle.

\section{Limiti}
I principali limiti sono:
\begin{itemize}
  \item la qualità del sistema dipende dalle regole di ancoraggio;
  \item la compressione può perdere informazioni rilevanti;
  \item due catene con uguale mutazione netta possono avere significato operativo diverso;
  \item il modello non sostituisce solutori numerici o modelli scientifici;
  \item la sintesi LLM delle policy richiede validazione umana o rule-based.
\end{itemize}
Il limite più critico è la perdita informativa. Per questo la firma compressa deve restare sempre collegata alla catena osservata e ai dati grezzi.

\section{Conclusione}
SemaGraph v0.6 formalizza un kernel deterministico per rappresentare catene di cambiamento in uno spazio finito di 64 stati booleani. La sua funzione non è predire, inferire o simulare direttamente la realtà, ma comprimere transizioni osservate in firme operative riproducibili. Questa struttura può ridurre inferenza ricorrente, migliorare auditabilità e abilitare policy synthesis controllata. La direzione Rust/WASM rende il nucleo compatibile con una futura esecuzione machine-near, mentre TypeScript resta il livello di specifica e integrazione.

\begin{thebibliography}{9}

\bibitem{shannon1948}
C. E. Shannon, ``A Mathematical Theory of Communication'', \emph{Bell System Technical Journal}, 1948.

\bibitem{hamming1950}
R. W. Hamming, ``Error Detecting and Error Correcting Codes'', \emph{Bell System Technical Journal}, 1950.

\bibitem{hopcroft2006}
J. E. Hopcroft, R. Motwani, J. D. Ullman, \emph{Introduction to Automata Theory, Languages, and Computation}, Pearson, 2006.

\bibitem{harel1987}
D. Harel, ``Statecharts: A Visual Formalism for Complex Systems'', \emph{Science of Computer Programming}, 1987.

\bibitem{alur1995}
R. Alur, D. L. Dill, ``A Theory of Timed Automata'', \emph{Theoretical Computer Science}, 1994.

\bibitem{russell2021}
S. Russell, P. Norvig, \emph{Artificial Intelligence: A Modern Approach}, Pearson, 2021.

\end{thebibliography}

\end{document}
